\section{Benchmark}
\label{sec:benchmark}

\begin{table*}[t]
\centering
\caption{Domains, change families and test questions. The no-effect column counts the test questions whose change, though valid, leaves the queried outcome unchanged.}
\label{tab:benchmark}
\footnotesize
\begin{tabular}{@{}l p{4.6cm} p{6.6cm} r r@{}}
\toprule
Domain & Source of inputs and rules & Change families (what a question changes) & Test qs. & No-effect \\
\midrule
Loan application & Published process model and simulation parameters~\cite{chapela2023zenodo}; 17 activities, loops, 6 worker pools & activity duration; timer delay; the route one case takes; routing probabilities; worker availability (workers added or unavailable); case removal & 120 & 24 \\
Pharmacy & Same source; 12 activities, loops, 3 worker pools & the same six families & 120 & 24 \\
Manufacturing & SupplyGraph production and orders~\cite{wasi2024supplygraph}; written dispatch policy~\cite{otftools2026} & order removal; production quantity; production timing; dispatch capacity & 96 & 16 \\
IT operations & Google 2011 cluster trace~\cite{reiss2011google}; written priority scheduler on 8 machines~\cite{otftools2026} & job removal; machine-failure removal; task priority; task resource request; machine removal & 120 & 20 \\
\midrule
Policy variants & Loan and Pharmacy with one more worker in one pool (P) or with the rules reworded (H) & the six workflow families & 91 & 22 \\
\bottomrule
\end{tabular}
\end{table*}


Evaluating whether a released answer is accurate required what-if questions with known ground-truth answers, posed to programs generated without seeing those questions. This section describes our evaluation domains, our methodology for guaranteeing and verifying exact ground-truth answers, and the candidate programs we assess.

\paragraph{Domains.} Table~\ref{tab:benchmark} lists the four domains. The two workflows come from an open repository of parameterized business-process models~\cite{chapela2023zenodo,chapela2024extraneous}; Manufacturing and IT operations reuse the documents of a prior benchmark of process records~\cite{otftools2026}, in which real production, order and cluster-trace logs run under written dispatch and schedule policies. A workflow run contains 15 cases, a manufacturing run covers one plant over 14 days, and an IT operations run captures one window of cluster activity.

\paragraph{Exact answers.} A what-if question has an exact answer only if the random parts of a run are pinned down, and a shared random seed does not do this. In a simulator that draws from one stream in execution order, such as Prosimos~\cite{lopezpintado2023prosimos}, adding one case shifts the draws of every later case, so the ``what if'' compares two unrelated runs. We therefore wrote the workflow semantics down first and built a deterministic reference engine for them. Every random quantity is tied to a stable identity (case, activity or decision, and visit) and printed in the document: planned work, planned waits, and the draw that decided each route. Removing a case leaves every other case's draws unchanged, and a new routing probability is applied to the same draw. This is the common-random-numbers technique~\cite{law2015simulation}: the question queries the same realized run under one change, and the document contains everything needed for an exact answer. Manufacturing and IT are deterministic by construction.

\paragraph{Checked reference answers.} Reference answers are computed by deterministic reference code, which was itself checked. For the workflows, a separate checker was written from the semantics specification alone, without the engine's code or outputs; engine and checker agree on 18 of 18 hand-worked cases, on 566 fixed and random scenarios in every output field, and on all 240 workflow test questions. Independent LLM agents given only the document text also matched 24 of 24 sampled answers. For manufacturing and IT, the original simulators and a second implementation fed only by values parsed back from the documents agree on all 216 test questions, and blind samples matched 16 of 16 (by hand) and 20 of 20 (by agents writing their own schedulers).

\paragraph{Questions and splits.} Each question is written from a fixed template for its family and selected by a fixed rule, without regard to how any system answers it. Data are split by run: workflow runs 1--4 are for development and runs 5--12 for testing, and IT windows that share jobs, or manufacturing periods that share order days, stay on the same side. The development questions (112 workflow, 108 manufacturing and IT) were used to design Vouch and are not reported as evidence.

\paragraph{Policy variants.} Programs are meant to be reused after the rules change, so the policy set asks questions under rules the program was not built for. In variant P (a policy change) one pool gets one more worker (Loan appraisers 2$\to$3; Pharmacy pharmacists 1$\to$2), stated in the document. In variant H (a harmless rewording) the parameters are unchanged and 24 sentences of the rules are reworded; four independent LLM agents reproduced every H run from the H text alone. A program that hard-codes the old pool size should fail on P, and nothing should change on H.

\paragraph{Platforms and programs.} Three platforms build programs. \otf~\cite{otftools2026} (first platform) compiles a document into functions that read its records and a simulator that runs them, and accepts a program if it reproduces hidden recorded outcomes closely enough. Its original code often failed on our documents for mechanical reasons, such as reading records from the wrong place, so we use a revised generation path, which we call R1, that makes record locations correct by construction and retries reading; R1 gives no hints about process logic and keeps the acceptance check. The unchanged original \otf accepted 32 of 60 builds; its programs answer only through \otf's own path and are reported in the artifact. \devs~\cite{chen2026devs} (second platform) generates discrete-event models in stages from a specification, to which we add one generic input format for records, plans and probabilities. The coding agent (third platform) is a single model asked to write a simulator for the same specification and run it once on a published example input. Claude Sonnet~5.5 and Haiku~5.5 build on every platform, and three open-weight models, gpt-oss-120b~\cite{openai2025gptoss}, Gemma~4 31B and Qwen3.5~9B, build on \otf only. Of 101 build attempts, 86 were accepted: 48 on \otf R1, where Qwen produced only 3 of 12, 22 on \devs and all 16 from the coding agent. Builds are independent and none was selected by its results. Builders had no access to test documents, reference answers or reference code, which we checked with sandbox isolation tests and by scanning every program for test-only names and values. The main test set has 9{,}792 program--question pairs from these 86 programs, and the policy set has 1{,}411 pairs from the 31 Claude-built workflow programs.
